\documentclass[11pt,aspectratio=169]{ctexbeamer}

\usetheme{default}
\setbeamertemplate{navigation symbols}{}

\usepackage{amsmath}
\usepackage{booktabs}

\definecolor{DeepBlue}{RGB}{25,74,122}
\definecolor{SoftInk}{RGB}{53,59,66}
\definecolor{PaperBg}{RGB}{249,248,245}
\definecolor{PanelBlue}{RGB}{236,243,249}

\setbeamercolor{normal text}{fg=SoftInk,bg=PaperBg}
\setbeamercolor{structure}{fg=DeepBlue}
\setbeamercolor{frametitle}{fg=DeepBlue,bg=PaperBg}
\setbeamercolor{title}{fg=DeepBlue}
\setbeamercolor{block title}{fg=DeepBlue,bg=PanelBlue}
\setbeamercolor{block body}{fg=SoftInk,bg=white}

\newcommand{\Core}[1]{
  \begin{block}{核心结论}
    #1
  \end{block}
}

\newcommand{\LectureMeta}[2]{
  \title{大模型推理基础设施}
  \subtitle{#1}
  \author{张书豪}
  \institute{华中科技大学计算机学院}
  \date{#2}
}


\LectureMeta{第 1 讲 课程导论}{研究生课程课件}

\begin{document}

\begin{frame}[plain]
  \titlepage
\end{frame}

\begin{frame}{教学目标}
  \Core{本讲建立课程全景视角，并说明大模型推理基础设施作为独立系统研究方向的形成背景。}
  \begin{itemize}
    \item 说明课程对象、主线和考核组织。
    \item 区分模型能力问题与系统能力问题。
    \item 建立“状态化推理系统”的统一理解。
    \item 给出整门课的统一问题地图。
  \end{itemize}
\end{frame}

\begin{frame}{为什么值得独立成课}
  \Core{大模型推理的关键矛盾已经从“模型能否生成”转向“服务能否稳定运行”。}
  \begin{itemize}
    \item 长上下文、MoE、RAG、智能体工作流重塑了推理负载。
    \item TTFT、P99、goodput、成本和稳定性成为核心系统指标。
    \item 真实收益来自全链路组织，而不是某个局部算子单点提速。
  \end{itemize}
\end{frame}

\begin{frame}{从模型问题到系统问题}
  \Core{当模型进入真实服务后，关注点会从“生成是否正确”转向“系统是否稳定兑现服务质量”。}
  \begin{itemize}
    \item 训练更关心收敛与效果。
    \item 推理更关心排队、缓存、执行路径与回归稳定性。
    \item 同一模型在不同运行时上会表现出完全不同的系统行为。
  \end{itemize}
\end{frame}

\begin{frame}{课程边界}
  \Core{本课程关注在线推理服务系统，而不是训练系统或纯算法优化。}
  \begin{itemize}
    \item 讨论请求如何进入系统、被调度、被执行和被回收。
    \item 讨论 KV Cache、知识缓存、会话记忆等状态对象。
    \item 讨论执行路径、异构平台和实验验证方法。
  \end{itemize}
\end{frame}

\begin{frame}{课程不重点展开什么}
  \Core{课程边界必须足够清楚，否则系统问题会被训练、提示词或纯算法问题淹没。}
  \begin{itemize}
    \item 不系统展开预训练、RLHF 和参数高效微调。
    \item 不把 prompt engineering 作为主体内容。
    \item 不把单个 kernel 调优等同于整门系统课程。
  \end{itemize}
\end{frame}

\begin{frame}{统一视角：状态化系统}
  \Core{推理系统的主要瓶颈往往来自共享状态的访问、保留、搬运和复用。}
  \begin{itemize}
    \item 模型参数之外，还有 KV Cache、知识缓存、会话记忆与调度元数据。
    \item 状态跨时间、跨请求、跨设备持续传播。
    \item 调度、缓存与执行路径因此连续耦合。
  \end{itemize}
\end{frame}

\begin{frame}{什么叫状态化系统}
  \Core{所谓状态化系统，是指系统当前行为持续受到历史对象与共享对象的影响。}
  \begin{itemize}
    \item KV 决定 decode 后续路径。
    \item 会话历史决定 prompt 长度和资源需求。
    \item 运行时统计信息反过来影响调度和资源治理。
  \end{itemize}
\end{frame}

\begin{frame}{课程主线}
  \Core{课程沿着工作负载、调度、状态管理、执行路径和验证方法五条线展开。}
  \begin{enumerate}
    \item 先建立工作负载和评价指标视角
    \item 再分解一次请求的完整生命周期
    \item 再进入调度、KV 与状态管理主问题
    \item 然后讨论异构执行、平台适配与实验方法
    \item 最后落到开源系统实践和课程项目
  \end{enumerate}
\end{frame}

\begin{frame}{课程问题地图}
  \Core{整门课都围绕一个总问题：如何在动态负载、持续状态和异构硬件约束下构建稳定可验证的推理服务。}
  \begin{itemize}
    \item 工作负载定义问题边界。
    \item 调度定义系统控制点。
    \item 状态组织定义容量与复用空间。
    \item 执行路径决定收益能否兑现。
  \end{itemize}
\end{frame}

\begin{frame}{课程组织方式}
  \Core{课程采用机制讲解、论文精读、代码实践与项目训练相结合的组织方式。}
  \begin{itemize}
    \item 平时阅读与课堂讨论
    \item 论文精读报告
    \item 最小复现实验或代码导读
    \item 课程项目中期与最终汇报
  \end{itemize}
\end{frame}

\begin{frame}{四层材料如何配合}
  \Core{整门课按“讲义定问题、tutorial 定路径、实验定证据、代码定控制点”的闭环推进。}
  \begin{itemize}
    \item 讲义与课件先建立工作负载、调度、状态和执行路径四条主线。
    \item Tutorials 把同一个问题落到 nano-vLLM-HUST 的最小代码对象和日志观察点。
    \item 课程实验要求学生围绕统一模板提交调用链、表格、日志和可否定解释。
    \item 教学代码同时提供参考实现、student\_exercises、assignment 包和 TA 汇总脚本，避免讲一套、做一套、交一套。
  \end{itemize}
\end{frame}

\begin{frame}{课程产出要求}
  \Core{课程产出应体现为可复查、可复用、可继续推进的研究材料。}
  \begin{itemize}
    \item 一份系统论文精读报告。
    \item 一份代码导读或最小复现实验记录。
    \item 一个课程项目及最终汇报材料。
  \end{itemize}
\end{frame}

\begin{frame}{讲解：这门课真正分析的对象是什么}
  \Core{《大模型推理引擎入门》不是一门“模型结构课”，而是一门“在线推理系统课”。}
  \begin{itemize}
    \item 我们关心的是请求怎样进入系统、怎样被调度、怎样携带状态、怎样形成用户看到的时延与吞吐。
    \item 同一个模型参数，在不同 runtime、不同调度器、不同 KV 管理策略下，会表现出完全不同的服务品质。
    \item 因此课程的基本分析单位不是某一层网络，而是“请求对象 + 控制点 + 状态对象 + 结果指标”。
  \end{itemize}
\end{frame}

\begin{frame}{讲解：导论课应带走的统一判断框架}
  \Core{从第一讲开始就要养成同一套判断顺序，否则后面会把现象、机制和结论混在一起。}
  \begin{enumerate}
    \item 先问：系统里长期存在的对象是什么，例如请求、KV、队列、会话状态。
    \item 再问：谁真正改变了系统行为，例如 scheduler、block manager、executor。
    \item 再问：收益如何传递到端到端指标，例如 TTFT、P99、goodput。
    \item 最后问：哪些日志、实验和代码证据能支持这条因果链。
  \end{enumerate}
\end{frame}

\begin{frame}{例子：为什么传统部署经验在 LLM 时代失效（1）}
  \Core{从分类模型到生成模型，服务系统的主矛盾从“单次前向吞吐”转向“跨轮状态组织”。}
  \begin{block}{阶段一：传统 CV / NLP 服务}
    单请求通常是一次前向，输入输出短，显存主要被模型参数占据。
  \end{block}
\end{frame}

\begin{frame}{例子：为什么传统部署经验在 LLM 时代失效（2）}
  \Core{从分类模型到生成模型，服务系统的主矛盾从“单次前向吞吐”转向“跨轮状态组织”。}
  \begin{block}{阶段二：早期小模型对话服务}
    即便开始有生成过程，请求状态较小，系统还能把问题近似看成短任务排队。
  \end{block}
\end{frame}

\begin{frame}{例子：为什么传统部署经验在 LLM 时代失效（3）}
  \Core{从分类模型到生成模型，服务系统的主矛盾从“单次前向吞吐”转向“跨轮状态组织”。}
  \begin{block}{阶段三：大模型在线服务}
    长上下文、连续 decode、KV 驻留和动态到达同时出现，服务问题不再是简单部署问题。
  \end{block}
\end{frame}

\begin{frame}{例子：真实服务里最先撞上的常见问题（1）}
  \Core{很多系统失效并不是“模型不会生成”，而是系统无法稳定承载状态和并发。}
  \begin{block}{现象一：晚高峰突然超时}
    请求并发上升后，TTFT 和 P99 同时恶化，用户感知是“模型卡住了”。
  \end{block}
\end{frame}

\begin{frame}{例子：真实服务里最先撞上的常见问题（2）}
  \Core{很多系统失效并不是“模型不会生成”，而是系统无法稳定承载状态和并发。}
  \begin{block}{现象二：显存没有满却接不了更多请求}
    连续空间、碎片、预留和块管理会让系统在理论容量之前就先失去弹性。
  \end{block}
\end{frame}

\begin{frame}{例子：真实服务里最先撞上的常见问题（3）}
  \Core{很多系统失效并不是“模型不会生成”，而是系统无法稳定承载状态和并发。}
  \begin{block}{现象三：换一套 runtime 行为完全不同}
    同一模型在不同服务栈上的表现差异，说明问题核心在系统组织而不在模型参数本身。
  \end{block}
\end{frame}

\begin{frame}{代码例子：最小请求从哪里进入系统（1）}
  \Core{课程不从复杂优化开始，而是先从一条最短代码路径建立控制点意识。}
  \begin{block}{入口代码}
    {\ttfamily\small
    \detokenize{from nanovllm import LLM, SamplingParams}\par
    \detokenize{llm = LLM(model_path, enforce_eager=True)}\par
    \detokenize{outputs = llm.generate(prompts, sampling_params)}\par
    }
  \end{block}
\end{frame}

\begin{frame}{代码例子：最小请求从哪里进入系统（2）}
  \Core{课程不从复杂优化开始，而是先从一条最短代码路径建立控制点意识。}
  \begin{block}{课堂应追问}
    这里的 \texttt{LLM.generate} 是最终控制点，还是对更深层 engine 的封装入口？
  \end{block}
\end{frame}

\begin{frame}{代码例子：最小请求从哪里进入系统（3）}
  \Core{课程不从复杂优化开始，而是先从一条最短代码路径建立控制点意识。}
  \begin{block}{系统含义}
    只要这一步还说不清，后续调度、KV、实验设计都会失去锚点。
  \end{block}
\end{frame}

\begin{frame}{代码例子：最小日志为什么足够有用（1）}
  \Core{导论课最有价值的代码例子不是复杂优化，而是能证明主路径存在的最小插桩。}
  \begin{block}{最小日志样式}
    {\ttfamily\small
    \detokenize{print(request_id, prompt_len, stage, timestamp)}\par
    \detokenize{print(batch_size, running_size, finished_count)}\par
    }
  \end{block}
\end{frame}

\begin{frame}{代码例子：最小日志为什么足够有用（2）}
  \Core{导论课最有价值的代码例子不是复杂优化，而是能证明主路径存在的最小插桩。}
  \begin{block}{为什么不一开始就打满日志}
    如果日志过多，学生会看到很多现象，但抓不住“谁拥有控制权”这个主问题。
  \end{block}
\end{frame}

\begin{frame}{代码例子：最小日志为什么足够有用（3）}
  \Core{导论课最有价值的代码例子不是复杂优化，而是能证明主路径存在的最小插桩。}
  \begin{block}{讲授重点}
    先用少量证据锁定路径，再逐步增加探针，这是系统调试和研究的共同方法。
  \end{block}
\end{frame}

\begin{frame}{反例：把系统问题误写成模型问题会怎样（1）}
  \Core{导论课最需要纠正的，就是把一切现象都归因于模型大小或算子快慢。}
  \begin{block}{误读一：输出慢是模型不够强}
    实际上，很多慢来自排队、prefill 压力和状态驻留，而不是模型回答质量。
  \end{block}
\end{frame}

\begin{frame}{反例：把系统问题误写成模型问题会怎样（2）}
  \Core{导论课最需要纠正的，就是把一切现象都归因于模型大小或算子快慢。}
  \begin{block}{误读二：换更快 kernel 就能解决服务问题}
    如果请求组织和状态管理失衡，再快的单点算子也难兑现到端到端收益。
  \end{block}
\end{frame}

\begin{frame}{反例：把系统问题误写成模型问题会怎样（3）}
  \Core{导论课最需要纠正的，就是把一切现象都归因于模型大小或算子快慢。}
  \begin{block}{导论阶段的任务}
    先把问题放回系统边界中，再谈优化对象和收益链路。
  \end{block}
\end{frame}

\begin{frame}{课堂练习：把抽象课程问题落到具体对象（1）}
  \Core{学生需要尽早学会把“大问题”翻译成“对象、控制点、证据”三件事。}
  \begin{block}{练习一}
    给出一个你见过的“模型卡顿”现象，判断它更像是排队、状态还是执行路径问题。
  \end{block}
\end{frame}

\begin{frame}{课堂练习：把抽象课程问题落到具体对象（2）}
  \Core{学生需要尽早学会把“大问题”翻译成“对象、控制点、证据”三件事。}
  \begin{block}{练习二}
    在 \texttt{example.py} 到 \texttt{llm\_engine.py} 之间，指出最值得先读的两个函数。
  \end{block}
\end{frame}

\begin{frame}{课堂练习：把抽象课程问题落到具体对象（3）}
  \Core{学生需要尽早学会把“大问题”翻译成“对象、控制点、证据”三件事。}
  \begin{block}{练习三}
    用一句话回答：为什么这门课要从“系统边界”而不是从“某篇论文机制”讲起？
  \end{block}
\end{frame}


\begin{frame}{课堂讨论}
  \Core{课程导论阶段的核心任务，是明确推理基础设施研究中的系统问题边界。}
  \begin{itemize}
    \item 为什么大模型服务问题不能简化为“把模型部署出去”？
    \item 用户感知到的“卡顿”背后通常是哪些系统现象？
    \item 哪些状态对象最容易成为瓶颈放大器？
  \end{itemize}
\end{frame}

\begin{frame}{本讲总结}
  \Core{将推理系统理解为状态化在线服务系统，有助于统一后续的调度、KV、记忆与异构执行分析。}
  \begin{itemize}
    \item 研究对象是完整服务链路而非单一模型。
    \item 研究目标是稳定、可验证的系统收益。
    \item 后续所有讲次都会回到“状态对象 + 控制点 + 收益链路”这个框架。
  \end{itemize}
\end{frame}

\end{document}